\section{基于预解式的谱距离估计}
\label{sec:error}

\pending{2026-11-06 重组：单一界面缺陷 $\varepsilon_{\partial,\ell_*}$ 改为分组缺陷 $\varepsilon_g$；迹估计的几何与证明移入附录~\ref{app:trace}，计入舍入的可靠计算移入附录~\ref{app:reliable}}
本节以第~\ref{sec:discretization}~节交付的候选为输入，估计
\begin{align}
  \mu&:=\widehat k^2
  \label{eq:candidate-mu}
\end{align}
到 $\sigma(A(\beta))$ 的距离。论证分三步：
\begin{enumerate}
  \item 用候选系数和实际使用的有限核表示在整个条带 $B$ 上构造分片场 $\widehat u$。它在
    每个均匀材料分片内精确满足 Helmholtz 方程，但在材料界面、人工竖直界面和条带上下
    接缝上留有小的跳跃，因而 $\widehat u\notin D(A(\beta))$
    （第~\ref{subsec:raw-field}~小节）。\pending{接缝跳跃为新增（2026-11-04）}
  \item 把界面分成有限个组，每组的跳跃汇总为界面缺陷 $\varepsilon_g$，并记
    $U:=\lVert\widehat u\rVert_{\mathcal H}$（第~\ref{subsec:jumps}~小节）。
  \item 借助正移位预解式的分组迹估计和 Green 恒等式证明
    $\operatorname{dist}(\mu,\sigma(A(\beta)))\leq2\mu Q/(U-Q)$，其中
    $Q=\sum_gC_g\varepsilon_g$，$C_g$ 只依赖几何、材料系数与 $\mu$
    （第~\ref{subsec:trace}、\ref{subsec:spectral-distance}~小节）。
\end{enumerate}
定理只使用 $\widehat u$ 的几条性质：它不要求离散方程有精确零解，不要求 $\widehat u$ 是
某个问题的唯一解，也不需要求解预解式方程。候选算法的全部细节只通过 $\varepsilon_g$ 和
$U$ 进入结论。第~\ref{subsec:computation}~小节说明这两个量的计算，
第~\ref{subsec:certified}~小节汇总可靠上界的输入。迹估计的几何与证明见
附录~\ref{app:trace}；计入舍入的可靠计算见附录~\ref{app:reliable}。



% ====================================================================
\subsection{原始场}
\label{subsec:raw-field}

\paragraph{定义参数与评价参数.}
本节始终区分两类参数。\emph{定义参数}是 Nystr\"om 网格 $h$（即密度样本数 $N_\eta$）、
Rayleigh 截断 $M$ 以及候选系数~\eqref{eq:candidate-vector}；它们一经固定，$\widehat u$
就是一个确定的函数。\emph{评价参数}是材料界面上的目标求积节点数 $N_q$、人工界面上的
模态截断 $M_{\mathrm{ev}}\geq M$、显式累加的胞元数 $N_{\mathrm{cell}}$ 以及核函数的
评价精度；它们只决定 $\varepsilon_g$ 和 $U$ 算得多准，不改变
$\widehat u$。候选向量的复数倍固定一次，后面所有量都由同一个场计算；由于
$\varepsilon_g$ 与 $U$ 都关于这个倍数一次齐次，最终的界与归一化无关。

\paragraph{连续密度.}
对每条材料曲线，候选计算在均匀网格 $t_r=r/N_\eta$ 上给出偶数个密度样本
$\{\eta_r\}_{r=0}^{N_\eta-1}$（$\tau$ 与 $-\omega$ 两个分量分别处理）。由离散 Fourier
变换定义
\begin{align}
  \eta_m
  &:=\frac1{N_\eta}\sum_{r=0}^{N_\eta-1}\eta_r
    \exp\!\left(-\frac{2\pi\ii mr}{N_\eta}\right),
  \qquad
  -\frac{N_\eta}{2}\leq m\leq\frac{N_\eta}{2}-1,
  \label{eq:exact-density-dft}
\end{align}
并令 $\eta(t):=\sum_m\eta_m\ee^{2\pi\ii mt}$，它在网格点上插值原样本。$\widehat u$ 中的
层势一律理解为这一三角多项式密度与精确核的连续积分。

不能直接把 Nystr\"om 求积和当作场：形如 $\sum_rw_rG(\mathbf z,\mathbf z_r)\omega_r$ 的
离散和虽然逐片满足 Helmholtz 方程，却在每个节点处带有点源和偶极子奇性。
\pending{说法更正（2026-11-04）}点源的对数奇性仍局部平方可积，偶极子的 $1/r$ 奇性
一般不平方可积；两者都不属于直到所在界面的 $H^2$，下面的定理无从应用。连续密度给出的场在界面两侧具有正常的单侧迹，
由跳跃关系~\eqref{eq:interface-jump-operator} 给出；其中的局部项要求在任意界面点取得
密度值，三角插值正好提供这些值。

\paragraph{中心胞元中的场.}
在 $\mathcal C_0$ 中，$\widehat u$ 由式~\eqref{eq:cell-layer-representation}
（$\Omega=\Omega_0$，$u_{\mathrm b}=u_{\mathrm{hom}}$）定义，其中 $\tau,\omega$ 为上述
连续密度，$u_{\mathrm{hom}}$ 只含 $|m|\leq M$ 的项、振幅为 $a^-_0,b^+_0$，全部波数用
$k=\widehat k$ 计算。

\paragraph{实际核表示.}
\pending{场的定义改为实际的增广 MFS 表示（2026-11-04 新增）}
实现中背景区域的准周期核由增广 MFS 近似（第~\ref{subsec:layer-potentials}~小节），
因此本文把实际使用的有限表示本身定义为场。在每个参考胞元 $\mathcal C$ 的背景区域中令
\begin{align}
  \widehat u
  &=u_{\mathrm{near}}
    +\sum_{q=1}^{N_p}\alpha_q\,G^{(\kappa_{\mathrm e})}(\cdot,\mathbf p_q)
    +\sum_{q=1}^{N_w}b_q\,\ee^{\ii\boldsymbol\zeta_q\cdot\mathbf z}
    +u_{\mathrm b},
  \qquad\kappa_{\mathrm e}=\widehat k\sqrt{\rho_{\mathrm e}},
  \label{eq:mfs-field}
\end{align}
其中 $u_{\mathrm{near}}$ 是本胞元曲线及其有限个近邻 $y$ 周期像上的\emph{连续}单层与
双层势，$\mathbf p_q$ 是代理点，$\boldsymbol\zeta_q\cdot\boldsymbol\zeta_q=\kappa_{\mathrm e}^2$，
$\alpha_q,b_q$ 是固定的有限复系数，$u_{\mathrm b}$ 是有限 Rayleigh 场；介质柱内仍是
波数 $\kappa_{\mathrm c}$ 的连续层势。系数线性依赖于胞元的端口状态，左右两侧各用自己的
几何、$L_\pm$ 与 $\mathsf P_\pm$。要求：(i) 本胞元以外的源曲线和全部代理点与该胞元的
闭包有正距离；(ii) 所有背景项都用同一波数 $\kappa_{\mathrm e}$。保存的浮点系数视为精确
的有理数；最小二乘拟合、QZ 分解与求解器是否准确不是定理的前提，其效果全部由该场的界面
跳跃衡量。若把代理源系数与密度、求积权的乘积先求和得到 $\alpha_q$，所得场一般与
“连续源依赖的 MFS 系数再精确积分”的场不同，界面跳跃必须对同一个场计算。

\paragraph{本文采用的表示：胞元核与按场拟合.}
\pending{具体选择（2026-11-05 新增）}记 $q:=\ee^{\ii\beta d}$。式~\eqref{eq:mfs-field} 中取
\begin{align}
  u_{\mathrm{near}}
  &=\sum_{m=-1}^{1}q^m\bigl(\mathcal S^{(\kappa_{\mathrm e})}[\omega]
    +\mathcal D^{(\kappa_{\mathrm e})}[\tau]\bigr)(\cdot-m\,d\,e_y),
  \label{eq:cell-kernel-near}
\end{align}
即本胞元曲线及其上下两个周期像上的连续层势，权重为准周期相位；不作源点与目标点之差的
回卷，因此胞元内没有切换线。代理点 $\mathbf p_q$（$N_p=480$）固定在矩形
$\{|x|=1.45\}\cup\{|y|=1.35\}$ 上，每边 $120$ 个；不加平面波（$N_w=0$）。在当前几何中，
上下两个像与闭胞元的距离至少 $0.24$，代理点至少 $0.85$，满足条件 (i)。

系数 $\alpha=(\alpha_q)$ 对整个场拟合，而不是对核拟合：在接缝所在直线 $y=\pm d/2$ 的
线段 $|x|\leq1.1$ 上取 $240$ 对配点，要求
$u_{\mathrm{near}}+\sum_q\alpha_qG^{(\kappa_{\mathrm e})}(\cdot,\mathbf p_q)$ 的值与
$\partial_y$ 满足准周期关系；在辅助直线 $x=\pm1.1$ 上再各取 $240$ 个配点，要求散射部分与
出射 Rayleigh 波（$|m|\leq32$）匹配，以排除胞元内的齐次解。由
式~\eqref{eq:cell-kernel-near} 的相消，接缝条件的右端只含 $q^{-1}$ 与 $q^2$ 两项像。辅助
出射振幅只是拟合中的未知量，不进入场。所得 $1440\times610$ 系统按列缩放后用相对截断
$10^{-14}$ 的奇异值分解求解。系统矩阵只依赖 $\kappa_{\mathrm e}$、$\beta$ 与几何，因此一次
分解同时处理中心胞元的一个右端与两侧引线各 $97$ 个状态列：中心取 $\alpha=\alpha_0$，第
$j$ 个引线胞元取 $\alpha=\mathsf C_\pm\mathsf P_\pm^jc_\pm$，$\mathsf C_\pm\in\CC^{480\times97}$
固定，满足“系数线性依赖于端口状态”的要求。

这种表示即文献~\cite{BarnettGreengard2010} 的“近像层势加辅助源”结构。该文把辅助系数与
密度联立求解；本文的密度沿用第~\ref{sec:discretization}~节的候选，只在事后拟合
$\alpha$。按命题~\ref{prop:piecewise-h2} 与第~\ref{subsec:spectral-distance}~小节的定理，
场无须满足任何离散方程，拟合与密度之间的不一致全部反映在界面跳跃中。若改用求解时的
回卷核作为场，胞元内的切换线会带来一类额外的体缺陷，不在本节定理的覆盖范围内。

\paragraph{周期胞元中的场.}
对右侧基本胞元，给定入射振幅 $(a_{\mathrm L},b_{\mathrm R})$，取第~\ref{subsec:lead-cells}~
小节中的密度 $\mathsf B_+(a_{\mathrm L};b_{\mathrm R})$，按式~\eqref{eq:exact-density-dft}
化为连续密度，并以式~\eqref{eq:cell-layer-representation}（$u_{\mathrm b}$ 为
式~\eqref{eq:cell-incident}）定义胞元场 $u_+^{\mathrm{cell}}[a_{\mathrm L},b_{\mathrm R}]$；
它线性依赖于入射振幅。左侧同样定义 $u_-^{\mathrm{cell}}$。对 $v\in\CC^K$，把
$\mathsf Z_\pm v$ 写成 $(a_\pm(v);b_\pm(v))$，并令
\begin{align}
  u_+[v]&:=u_+^{\mathrm{cell}}\bigl[a_+(v),\,b_+(\mathsf P_+v)\bigr]\ \text{于 }\mathcal C_+,
  &u_-[v]&:=u_-^{\mathrm{cell}}\bigl[a_-(\mathsf P_-v),\,b_-(v)\bigr]\ \text{于 }\mathcal C_-.
  \label{eq:lead-cell-field}
\end{align}
在 $u_+[v]$ 中 $v$ 是左端口状态的坐标、$\mathsf P_+v$ 是右端口状态的坐标；$u_-[v]$ 中
$v$ 是右端口状态的坐标、$\mathsf P_-v$ 是左端口状态的坐标。由几何与端口锚定的平移
不变性，在 $\mathcal C_{j,\pm}$ 上令
\begin{align}
  \widehat u(\mathbf z\pm jL_\pm e_x)
  &:=u_\pm\bigl[\mathsf P_\pm^jc_\pm\bigr](\mathbf z),
  \qquad\mathbf z\in\mathcal C_\pm,\ j\geq0.
  \label{eq:half-guide-raw-field}
\end{align}

这一构造在 $|m|\leq M$ 的模态上自洽。以右侧为例，记
$z_j:=\mathsf Z_+\mathsf P_+^jc_+=(a_j;b_j)$。由式~\eqref{eq:propagation-plus}，
$\mathsf E_{\mathrm L}z_j=\mathsf E_{\mathrm R}z_{j+1}$，即
$b_j=R_{\mathrm L}a_j+T_{\mathrm{RL}}b_{j+1}$、
$a_{j+1}=T_{\mathrm{LR}}a_j+R_{\mathrm R}b_{j+1}$。与
式~\eqref{eq:cell-outgoing}--\eqref{eq:one-cell-scattering} 比较可知，第 $j$ 个胞元以
$(a_j,b_{j+1})$ 入射时，其出射振幅恰为 $b_j$ 与 $a_{j+1}$。因此共享端口两侧的场在
$|m|\leq M$ 上都具有振幅 $(a_{j+1},b_{j+1})$；而 $\mathsf M(\widehat k)$ 的端口行块使
中心场与 $z_0$ 在 $\Gamma_+$ 上近似匹配。剩余的不匹配来自 $|m|>M$ 的模态、连续密度与
离散密度的差别，以及 $\sigma_{\min}\ne0$。由于 $\mathsf P_\pm$ 的谱半径小于一，
$\widehat u$ 向两侧逐胞元几何衰减。

\paragraph{原始场的性质.}
称 $\widehat u$ 为\emph{原始场}。它满足：
\begin{enumerate}
  \item[(P1)] 在每个均匀材料分片内，
    \begin{align}
      (-\Delta-\mu\rho)\widehat u&=0,
      \label{eq:raw-piecewise-pde}
    \end{align}
    因为所用的核与 Rayleigh 波都是精确的 Helmholtz 解；
  \item[(P2)] \pending{原稿要求精确准周期，2026-11-04 删去}不要求 $\beta$-准周期条件
    精确成立；上下接缝上的失配作为界面跳跃计入 $\varepsilon_g$
    （第~\ref{subsec:jumps}~小节）；
  \item[(P3)] 在每个分片内直至其边界局部属于 $H^2$（命题~\ref{prop:piecewise-h2}），
    从而逐片 Green 恒等式与普通法向迹有意义；
  \item[(P4)] $\widehat u\in\mathcal H$，这由传播衰减保证，并在
    第~\ref{subsec:mass}~小节核验。
\end{enumerate}
但 $\widehat u\notin H^1_\beta(B)$，因为界面上的值跳跃一般不为零。

\begin{proposition}[逐片 $H^2$ 正则性]
  \label{prop:piecewise-h2}
  \pending{新增（2026-11-04）}设每条材料界面是正则、嵌入的解析闭曲线，与胞元边界有
  正距离，密度是三角多项式，且表示~\eqref{eq:mfs-field} 满足条件 (i)(ii)。则
  $\widehat u$ 在每个材料分片内直至边界局部属于 $H^2$，并在人工界面与接缝的每一侧有
  普通的值迹与法向迹。
\end{proposition}

\begin{proof}
  远处的像、代理源和平面波在闭分片附近光滑，只需考虑一条曲线 $\Gamma$ 上的
  $v=\mathcal S^{(\kappa)}[\omega]+\mathcal D^{(\kappa)}[\tau]$。它局部平方可积，且在分布
  意义下 $(\Delta+\kappa^2)v=-\omega\delta_\Gamma+\operatorname{div}(\tau\nu\delta_\Gamma)$。
  在 $\Gamma$ 的法向坐标中取 $w$：内侧为零，外侧为 $(\tau(s)-t\omega(s))\chi(t)$，$\chi$
  为光滑截断。$w$ 的值跳跃为 $\tau$、法向跳跃为 $-\omega$，故
  $(\Delta+\kappa^2)w=F-\omega\delta_\Gamma+\operatorname{div}(\tau\nu\delta_\Gamma)$，
  $F\in L^2_{\mathrm{loc}}$。于是 $(\Delta+\kappa^2)(v-w)=-F$，由常系数椭圆方程的内部
  正则性 $v-w\in H^2_{\mathrm{loc}}$，而 $w$ 在每侧直至 $\Gamma$ 光滑。人工界面与接缝附近
  的单侧胞元公式可以向外光滑延拓，故不存在角点正则性问题。
\end{proof}

% ====================================================================
\subsection{界面跳跃与两个输入量}
\label{subsec:jumps}

\paragraph{分片与界面族.}
条带被分为下列有界均匀材料分片：$\mathcal C_0\setminus\overline{\Omega_0}$ 与 $\Omega_0$
（若非空），以及各 $\mathcal C_{j,\pm}\setminus\overline{\Omega_{j,\pm}}$ 与
$\Omega_{j,\pm}$。分片的边界由材料界面、人工竖直界面和条带上下边界的线段组成。
\pending{接缝作为第三类界面（2026-11-04 修改）}把条带看作准周期圆柱：上边界
$\Sigma^\uparrow$ 与下边界 $\Sigma^\downarrow$ 粘合成一条接缝，$y>d/2$ 一侧的场用
$\ee^{\ii\beta d}\widehat u(x,y-d)$ 表示。对每个胞元 $\mathcal C$（$\mathcal C_0$ 或
$\mathcal C_{j,\pm}$），记 $\Sigma_{\mathcal C}$ 为其 $x$ 区间上的接缝段。定义互不相交的
界面族
\begin{align}
  \mathscr E
  &:=\mathscr E_{\mathrm{mat}}\mathbin{\dot\cup}\mathscr E_{\mathrm{art}}
    \mathbin{\dot\cup}\mathscr E_\Sigma,
  &\mathscr E_{\mathrm{mat}}
  &:=\{\partial\Omega_0\}\cup\{\partial\Omega_{j,\pm}\}_{j\geq0},
  \notag\\
  \mathscr E_{\mathrm{art}}
  &:=\{\Gamma_-,\Gamma_+\}\cup\{\Gamma_{j,\pm}\}_{j\geq0},
  &\mathscr E_\Sigma
  &:=\{\Sigma_{\mathcal C_0}\}\cup\{\Sigma_{\mathcal C_{j,\pm}}\}_{j\geq0},
  \label{eq:interface-family}
\end{align}
$\Omega_0=\varnothing$ 时省略 $\partial\Omega_0$。每条共享曲线恰好计一次；每个接缝段的
上、下两边只配成一条接缝。若场精确准周期，接缝上的跳跃为零，就回到原稿的设定。

\paragraph{定向与跳跃.}
在每条 $\Gamma\in\mathscr E$ 上指定侧 $1$、侧 $2$，并取由侧 $1$ 指向侧 $2$ 的单位法向
$\nu_\Gamma$：材料界面上侧 $1$ 为介质柱内部；人工界面上侧 $1$ 为较靠近中心的胞元，
$\nu_\Gamma$ 指向相应无穷远（在 $\Gamma_\pm$ 上即 $\nu_\pm$）；接缝上侧 $1$ 为
$y<d/2$ 一侧，$\nu_\Sigma=e_y$。定义
\begin{align}
  \jump{\widehat u}_\Gamma
  &:=\widehat u^{(2)}|_\Gamma-\widehat u^{(1)}|_\Gamma,
  &\jump{\partial_\nu\widehat u}_\Gamma
  &:=\partial_{\nu_\Gamma}\widehat u^{(1)}|_\Gamma
    -\partial_{\nu_\Gamma}\widehat u^{(2)}|_\Gamma.
  \label{eq:raw-interface-jumps}
\end{align}
两者的符号约定不同，是为了使后面 Green 恒等式中的两项都带正号。在接缝上，记
$q:=\ee^{\ii\beta d}$，式~\eqref{eq:raw-interface-jumps} 即
\begin{align}
  \jump{\widehat u}_\Sigma(x)&=q\,\widehat u(x,-d/2)-\widehat u(x,d/2),
  &\jump{\partial_\nu\widehat u}_\Sigma(x)
  &=\partial_y\widehat u(x,d/2)-q\,\partial_y\widehat u(x,-d/2).
  \label{eq:seam-jumps}
\end{align}在材料界面上，
$\jump{\widehat u}_\Gamma$ 与 $-\jump{\partial_\nu\widehat u}_\Gamma$ 恰是
式~\eqref{eq:interface-jump-operator} 的两个分量（$\partial\Omega_0$ 上取
$u_{\mathrm b}=u_{\mathrm{hom}}$，周期胞元中取入射场），其中的算子是作用于连续密度的
连续边界积分算子。因此 $J(\widehat k)\xi\approx0$ 只表示跳跃在 Nystr\"om 节点上近似为零。
类似地，$\mathsf M(\widehat k)$ 的端口行块近似就是 $\Gamma_\pm$ 上 $|m|\leq M$ 模态的
跳跃系数。


\paragraph{分组与两个输入量.}
\pending{2026-11-06：原稿只用一个参考长度 $\ell_*$ 定义 $\varepsilon_{\partial,\ell_*}$，现改为分组定义；只取一组时即原定义}
值跳跃与法向导数跳跃的量纲相差一个长度。把 $\mathscr E$ 分成有限个互不相交的组
$\mathscr E_g$（本文取五组：中心材料曲线、左引线材料曲线、右引线材料曲线、人工竖直界面、
接缝），每组取一个参考长度 $\ell_g>0$，定义
\begin{align}
  \varepsilon_g^2
  &:=\sum_{\Gamma\in\mathscr E_g}
  \left(
    \lVert\jump{\partial_\nu\widehat u}_\Gamma\rVert_{L^2(\Gamma)}^2
    +\ell_g^{-2}\lVert\jump{\widehat u}_\Gamma\rVert_{L^2(\Gamma)}^2
  \right),
  &U&:=\lVert\widehat u\rVert_{\mathcal H}.
  \label{eq:boundary-defect-norm}
\end{align}
与矩阵残量相比，$\varepsilon_g$ 还计入了材料曲线上节点之间的跳跃、$|m|>M$ 的模态、
接缝上的准周期失配，以及半波导中全部材料、人工界面与接缝。小的 $\sigma_{\min}$ 只是
$\varepsilon_g$ 小的部分证据，后面的定理只使用后者。


% ====================================================================
\subsection{正移位预解式的分组迹估计}
\label{subsec:trace}

\pending{2026-11-06：本小节只保留分组命题的陈述；法向坐标、Jacobian、覆盖重数的图示与讨论以及证明移入附录~\ref{app:trace}；原全界面命题（原命题 5.2）是分组命题的特例，只在工作版附录中保留}
证明谱距离定理时，要用迹不等式把界面上的 $L^2$ 范数控制到条带上的体积分。对每组
界面，在每条界面两侧取半宽 $\delta_g$ 的“领子”，即沿法向伸出的带状邻域；在每条法向
线段上用加权一维迹不等式，再对整条界面积分。需要的几何条件是：
\begin{itemize}
  \item 法向映射 $\boldsymbol\Phi_\Gamma(s,t):=\boldsymbol r_\Gamma(s)+t\nu_\Gamma(s)$
    在 $|t|<\delta_g$ 上单射，即领子不自交（$\boldsymbol r_\Gamma$ 为弧长参数，$\nu_\Gamma$
    为第~\ref{subsec:jumps}~小节规定的法向）；
  \item 带符号曲率满足 $\lvert\kappa_\Gamma\rvert\leq\kappa_g$ 且 $h_g:=\delta_g\kappa_g<1$，
    从而领子内的 Jacobian $1-t\kappa_\Gamma(s)$ 为正；
  \item 组内覆盖重数至多 $N_g$，即任一点至多属于组内 $N_g$ 个领子。
\end{itemize}
人工界面是直线段，曲率为零，领子是竖直带；接缝的领子是准周期圆柱上的水平带。领子只
用于辅助场 $z\in H^2_\beta(B)$（式~\eqref{eq:operator-domain-h2}），不用于不连续的
$\widehat u$，所以领子可以穿过其他界面；不同组的领子重叠也不需要跨组的重数，因为各组
分别由同一个全局体范数控制。


\begin{proposition}[分组迹估计]
  \label{prop:grouped-trace-bound}
  把 $\mathscr E$ 分成有限个组 $\mathscr E_g$（例如中心材料曲线、左右引线材料曲线、人工
  竖直界面、接缝；同一条曲线也可按弧长分段）。对每组设：法向领子半宽为 $\delta_g$ 且
  法向映射单射，$\lvert\kappa_\Gamma\rvert\leq\kappa_g$，$h_g:=\delta_g\kappa_g<1$，组内
  覆盖重数至多 $N_g$。令
  \begin{align}
    \tau_g&:=\frac{\delta_g}{6}+\frac{\delta_gh_g^2}{40(1-h_g^2)},
    &D_{0,g}&:=\frac{N_g(1+\eta_g)}{2\delta_g},
    &D_{1,g}&:=N_g(1+\eta_g^{-1})\tau_g,
    \label{eq:jacobian-trace-constants}
  \end{align}
  其中 $\eta_g>0$ 任取。则对 $v\in D(A(\beta))$，
  \begin{align}
    \sum_{\Gamma\in\mathscr E_g}\lVert v\rVert_{L^2(\Gamma)}^2
    &\leq D_{0,g}\lVert v\rVert_{L^2(B)}^2+D_{1,g}\lVert\nabla v\rVert_{L^2(B)}^2,
    \notag\\
    \sum_{\Gamma\in\mathscr E_g}\lVert\partial_{\nu_\Gamma}v\rVert_{L^2(\Gamma)}^2
    &\leq D_{0,g}\lVert\nabla v\rVert_{L^2(B)}^2+D_{1,g}\lVert D^2v\rVert_{L^2(B)}^2 .
    \label{eq:jacobian-trace-bound}
  \end{align}
  令 $p_g(\lambda):=a_g+b_g\lambda+c_g\lambda^2$，$a_g:=D_{0,g}/\rin$，
  $b_g:=D_{1,g}+\ell_g^2D_{0,g}$，$c_g:=\ell_g^2D_{1,g}\rax$，以及
  \begin{align}
    C_g^2&:=\sup_{\lambda\geq\lambda_0}\frac{p_g(\lambda)}{(\lambda+\mu)^2},
    \label{eq:joint-multiplier-constant}
  \end{align}
  $\lambda_0$ 如式~\eqref{eq:spectral-bottom}（取 $\lambda_0=0$ 也成立）。则
  $T_g(v)^2:=\sum_{\Gamma\in\mathscr E_g}(\lVert v\rVert_{L^2(\Gamma)}^2
  +\ell_g^2\lVert\partial_{\nu_\Gamma}v\rVert_{L^2(\Gamma)}^2)
  \leq C_g^2\lVert(A(\beta)+\mu)v\rVert_{\mathcal H}^2$。
\end{proposition}

证明见附录~\ref{app:trace-proof}。式~\eqref{eq:joint-multiplier-constant} 的上确界不需要扫描：只需比较 $\lambda=\lambda_0$、
$\lambda\to\infty$（值 $c_g$）以及驻点 $\lambda_*=(2a_g-b_g\mu)/(2c_g\mu-b_g)$（若
$\lambda_*\geq\lambda_0$）。各组可以选不同的 $\ell_g$。

后文简记
\begin{align}
  Q&:=\sum_gC_g\varepsilon_g .
  \label{eq:trace-defect-product}
\end{align}
$C_g$ 只依赖 $\mu$、$\rin$、$\rax$、$\ell_g$ 与几何量，不是由候选残量拟合的量。

% ====================================================================
\subsection{谱距离定理}
\label{subsec:spectral-distance}

\begin{theorem}[原始界面缺陷的谱距离界]
  \label{thm:spectral-distance}
  \pending{2026-11-06：陈述与证明改为分组形式}设界面分组（含接缝）满足
  命题~\ref{prop:grouped-trace-bound} 的几何条件，$\mu>0$，$Q$ 如
  式~\eqref{eq:trace-defect-product}。设非零场 $\widehat u$ 满足 (P1)--(P4)，且每个
  $\varepsilon_g<\infty$。若
  \begin{align}
    Q&<U,
    \label{eq:positive-denominator}
  \end{align}
  则
  \begin{align}
    \operatorname{dist}\bigl(\mu,\sigma(A(\beta))\bigr)
    &\leq\frac{2\mu Q}{U-Q}.
    \label{eq:spectral-distance-bound}
  \end{align}
  这是到整个谱的距离上界，不需要预先知道本质谱或谱隙。
\end{theorem}

\begin{proof}
  若 $\mu\in\sigma(A(\beta))$，结论显然。以下设
  $\delta:=\operatorname{dist}(\mu,\sigma(A(\beta)))>0$，并只为证明引入辅助场
  \begin{align}
    z&:=(A(\beta)-\mu)^{-1}\widehat u\in D(A(\beta)),
    \label{eq:auxiliary-resolvent-field}
  \end{align}
  计算中不需要求解它。证明中简写 $u=\widehat u$。思路是：以 $z$ 为 Green 恒等式的测试
  场，把 $U^2$ 写成跳跃与 $z$ 的迹的配对，再用命题~\ref{prop:grouped-trace-bound}
  控制 $z$ 的迹。

  \smallskip\noindent\textbf{第一步：逐片恒等式。}
  在每个均匀分片内 $\Delta u=-\mu\rho u$，而由 $(A(\beta)-\mu)z=u$ 得
  $-\Delta z-\mu\rho z=\rho u$。$\mu,\rho$ 为实数，取共轭后乘以 $u$ 并代入前式，得
  $\rho|u|^2=\overline z\,\Delta u-u\,\Delta\overline z$。这里只用到 $u$ 的逐片方程。
  对有界分片 $D$ 及其外法向 $\nu_D$，Green 第二恒等式给出
  \begin{align*}
    \int_D\rho|u|^2
    =\int_{\partial D}\bigl(\overline z\,\partial_{\nu_D}u-u\,\partial_{\nu_D}\overline z\bigr)\dd s.
  \end{align*}

  \smallskip\noindent\textbf{第二步：合并共享界面。}
  在 $\Gamma\in\mathscr E$ 上，侧 $1$ 的外法向为 $\nu_\Gamma$，侧 $2$ 的为 $-\nu_\Gamma$。
  由引理~\ref{lem:global-h2}，$z\in H^2(B)$，其值迹与法向迹在两侧相同。两侧边界被积
  函数相加得
  \begin{align*}
    \bigl(\overline z\,\partial_{\nu_\Gamma}u^{(1)}-u^{(1)}\partial_{\nu_\Gamma}\overline z\bigr)
    -\bigl(\overline z\,\partial_{\nu_\Gamma}u^{(2)}-u^{(2)}\partial_{\nu_\Gamma}\overline z\bigr)
    =\jump{\partial_\nu u}_\Gamma\overline z+\jump{u}_\Gamma\partial_{\nu_\Gamma}\overline z.
  \end{align*}
  \pending{接缝项（2026-11-04 修改）}接缝也按这种方式合并：$z$ 精确准周期，按
  $y>d/2$ 一侧的准周期延拓，其值迹与法向迹在接缝两侧相同，而 $u$ 的两侧值由
  式~\eqref{eq:seam-jumps} 给出。直接代入上、下边界的外法向可验证，两条边界项之和恰为
  $\int_{\Sigma_{\mathcal C}}(\jump{\partial_\nu u}_\Sigma\overline z
  +\jump{u}_\Sigma\partial_y\overline z)\,\dd x$，与其他界面的形式相同。若 $u$ 精确
  准周期，这一项为零。

  \smallskip\noindent\textbf{第三步：截断。}
  取实值 $\chi_R\in C_c^\infty(\RR)$，$0\leq\chi_R\leq1$，在 $|x|\leq R$ 上为 $1$，在
  $|x|\geq2R$ 上为 $0$，$\lVert\chi_R'\rVert_\infty=O(R^{-1})$，
  $\lVert\chi_R''\rVert_\infty=O(R^{-2})$。把前两步中的 $z$ 换成 $\chi_Rz$，只需对有限
  个分片求和，最外侧竖直边界没有贡献；接缝上 $\partial_{\nu_\Sigma}\chi_R=\partial_y\chi_R=0$。由
  $\Delta(\chi_R\overline z)=\chi_R\Delta\overline z+2\chi_R'\partial_x\overline z
  +\chi_R''\overline z$ 及
  $\partial_{\nu_\Gamma}(\chi_R\overline z)=\chi_R\partial_{\nu_\Gamma}\overline z
  +(\partial_{\nu_\Gamma}\chi_R)\overline z$，
  \begin{align*}
    \int_B\chi_R\rho|u|^2
    ={}&\sum_{\Gamma\in\mathscr E}\int_\Gamma\chi_R\bigl(\jump{\partial_\nu u}_\Gamma\overline z
      +\jump{u}_\Gamma\partial_{\nu_\Gamma}\overline z\bigr)\dd s\\
    &+\underbrace{\int_Bu\bigl(2\chi_R'\partial_x\overline z+\chi_R''\overline z\bigr)}_{=:I_R}
    +\underbrace{\sum_{\Gamma\in\mathscr E}\int_\Gamma
      \jump{u}_\Gamma(\partial_{\nu_\Gamma}\chi_R)\overline z\,\dd s}_{=:J_R}.
  \end{align*}

  \smallskip\noindent\textbf{第四步：令 $R\to\infty$。}
  记 $T_g(z)$ 如命题~\ref{prop:grouped-trace-bound}。对每组用加权 Cauchy--Schwarz
  不等式（法向跳跃配 $z$，值跳跃乘 $\ell_g^{-1}$ 配 $\ell_g\partial_\nu z$），
  \begin{align}
    \sum_{\Gamma\in\mathscr E_g}\int_\Gamma
      \bigl(|\jump{\partial_\nu u}_\Gamma|\,|z|+|\jump{u}_\Gamma|\,|\partial_{\nu_\Gamma}z|\bigr)\dd s
    \leq\varepsilon_gT_g(z),
    \qquad T_g(z)\leq C_g\lVert(A(\beta)+\mu)z\rVert_{\mathcal H}<\infty.
    \label{eq:weighted-cauchy-schwarz}
  \end{align}
  由 $u\in L^2(B)$ 与 $z\in H^2(B)$，$I_R\to0$；同理
  $|J_R|\leq\lVert\chi_R'\rVert_\infty\sum_g\ell_g\varepsilon_gT_g(z)\to0$。这正是迹估计在
  去掉截断时的作用。由式~\eqref{eq:weighted-cauchy-schwarz}，主界面项可由支配收敛取
  极限；左端同理趋于 $U^2$。于是
  \begin{align}
    U^2&=\sum_{\Gamma\in\mathscr E}\int_\Gamma
      \bigl(\jump{\partial_\nu u}_\Gamma\overline z
        +\jump{u}_\Gamma\partial_{\nu_\Gamma}\overline z\bigr)\dd s.
    \label{eq:raw-green-identity}
  \end{align}

  \smallskip\noindent\textbf{第五步：预解式估计。}
  由谱定理 $\lVert(A(\beta)-\mu)^{-1}\rVert=\delta^{-1}$，故
  $\lVert z\rVert_{\mathcal H}\leq U/\delta$；又 $(A(\beta)+\mu)z=u+2\mu z$，所以
  $\lVert(A(\beta)+\mu)z\rVert_{\mathcal H}\leq(1+2\mu/\delta)U$。由
  式~\eqref{eq:raw-green-identity} 与 \eqref{eq:weighted-cauchy-schwarz}，
  \begin{align*}
    U^2\leq\sum_g\varepsilon_gT_g(z)
    \leq Q\Bigl(1+\frac{2\mu}{\delta}\Bigr)U.
  \end{align*}
  除以 $U$ 并利用 $U-Q>0$ 得 $\delta(U-Q)\leq2\mu Q$。
\end{proof}

界的右端关于 $\widehat u$ 的倍数不变；当所有 $\varepsilon_g\to0$ 时，它以 $2\mu Q/U$
线性趋于零。定理对每组 $\ell_g>0$ 都成立，$\ell_g$ 可用来平衡两类跳跃。


\begin{remark}[体缺陷]
  若使用不精确满足 (P1) 的分片场，令 $r=(-\Delta-\mu\rho)\widehat u$、
  $\varepsilon_{\mathrm{vol}}:=\lVert\rho^{-1}r\rVert_{\mathcal H}$，则第一步多出
  $\int_Br\overline z$，其绝对值不超过 $\varepsilon_{\mathrm{vol}}U/\delta$，同一证明给出
  $\operatorname{dist}(\mu,\sigma(A(\beta)))\leq(2\mu Q+\varepsilon_{\mathrm{vol}})/(U-Q)$。
  本文的原始场体缺陷为零。
\end{remark}

% ====================================================================
\subsection{界面缺陷与场范数的计算}
\label{subsec:computation}

$\varepsilon_g^2$ 是对无穷多条界面的求和。下面依次说明单条材料界面与单条人工界面上的
积分，再用传播矩阵把无穷求和化为有限计算加可控尾项。本小节给出的求积是数值估计；
可靠上界见第~\ref{subsec:reliable-evaluation}~小节与附录~\ref{app:reliable}。


\subsubsection{材料界面}

设 $\Gamma=\partial\Omega$ 参数化为 $\boldsymbol r(t)$，$t\in[0,1)$。在目标点
$\mathbf z_q=\boldsymbol r(t_q)$，$t_q=q/N_q$，处按式~\eqref{eq:interface-jump-operator}
计算两类量：
\begin{itemize}
  \item 局部项 $\tau(t_q)$、$\omega(t_q)$：由三角插值~\eqref{eq:exact-density-dft}
    直接求值；
  \item 积分项：取 $N_q\geq N_\eta$，把密度零填充插值到 $N_q$ 点网格（三角多项式次数
    小于 $N_\eta/2$，插值精确），再用 Kress 核分裂求积。由
    第~\ref{subsec:layer-potentials}~小节，四个算子差的核至多对数奇异，无须超奇异求积。
\end{itemize}
背景场 $u_{\mathrm b}$ 及其法向导数直接求值。得到 $\jump{\widehat u}_\Gamma(\mathbf z_q)$
与 $\jump{\partial_\nu\widehat u}_\Gamma(\mathbf z_q)$ 后，用周期梯形公式
\begin{align}
  \int_\Gamma\bigl(|\jump{\partial_\nu\widehat u}_\Gamma|^2
    +\ell_g^{-2}|\jump{\widehat u}_\Gamma|^2\bigr)\dd s
  \approx\frac1{N_q}\sum_{q=0}^{N_q-1}
  \bigl(|\jump{\partial_\nu\widehat u}_\Gamma(\mathbf z_q)|^2
    +\ell_g^{-2}|\jump{\widehat u}_\Gamma(\mathbf z_q)|^2\bigr)|\boldsymbol r'(t_q)|.
  \label{eq:material-jump-quadrature}
\end{align}
这里有两层求积：内层评价层势在目标点的值，外层计算 $L^2$ 范数。$N_q$ 可以远大于
$N_\eta$ 而不改变密度。

\subsubsection{人工界面与接缝}

\pending{2026-11-06：原模态 Parseval 计算与模态尾界只适用于精确准周期 Green 函数，对实际场不成立，移入工作版保留段}
人工界面与接缝都是直线段，两侧的场都可在段上任意点直接求值：本胞元的层势只在材料
曲线上奇异，近像、代理点源与 Rayleigh 波在段附近光滑。跳跃函数一般不满足端点准周期
匹配，其 Fourier 系数通常只有 $1/|m|$ 的尾，所以不按模态截断，而是像材料界面一样分成
面板，对跳跃函数作带误差界的分段求积（第~\ref{subsec:reliable-evaluation}~小节）。


\begin{workonly}
\pending{本小节的模态尾界对实际场不成立（2026-11-04）}下面的模态计算与尾界
式~\eqref{eq:artificial-tail-bound} 用到精确准周期 Green 函数的 Rayleigh 结构：层势的
法向系数等于值系数乘 $\pm\ii\gamma_m$，高阶模态按 $\ee^{-|\gamma_m|\delta_\Gamma}$ 衰减。
对增广 MFS 场这两点都不成立：跳跃函数一般不满足端点准周期匹配，其 Fourier 系数通常只有
$1/|m|$ 的尾。Parseval 等式~\eqref{eq:artificial-jump-parseval} 对实际跳跃仍成立，但
尾项须改用实际跳跃的端点差与导数界，或者直接在墙段上对跳跃函数作带误差界的分段求积。
本文的实现采用后者；以下原文保留，只适用于精确准周期场。

设 $\Gamma=\{x_\Gamma\}\times I_d$。每侧的场在 $\Gamma$ 附近由 Rayleigh 波与该侧胞元的
准周期层势组成，层势源全在 $\Gamma$ 的同一侧。按 $\psi_m$ 展开每侧场的值迹与
$\partial_x$ 迹：Rayleigh 波直接给出振幅；层势的振幅由
式~\eqref{eq:rayleigh-coefficient-right} 的连续形式给出（$x_p=x_\Gamma$），其值系数为
$\phi_m$，$\partial_x$ 系数为 $\pm\ii\gamma_m\phi_m$（源在左侧取正）。两侧相减并按
$\nu_\Gamma$ 取法向导数，得到 $\jump{\widehat u}_{\Gamma,m}$ 与
$\jump{\partial_\nu\widehat u}_{\Gamma,m}$。由于 $\{\psi_m\}$ 正交归一，
\begin{align}
  \int_\Gamma\bigl(|\jump{\partial_\nu\widehat u}_\Gamma|^2
    +\ell_*^{-2}|\jump{\widehat u}_\Gamma|^2\bigr)\dd s
  &=\sum_{m\in\ZZ}\bigl(|\jump{\partial_\nu\widehat u}_{\Gamma,m}|^2
    +\ell_*^{-2}|\jump{\widehat u}_{\Gamma,m}|^2\bigr).
  \label{eq:artificial-jump-parseval}
\end{align}

只显式计算 $|m|\leq M_{\mathrm{ev}}$ 的项，其余项不能直接设为零。取
$M_{\mathrm{ev}}\geq M$ 足够大，使 $|m|>M_{\mathrm{ev}}$ 时 $\beta_m^2>\kappa_{\mathrm e}^2$，
于是 $\gamma_m=\ii|\gamma_m|$，且这些模态只有层势贡献。记 $\delta_\Gamma>0$ 为 $\Gamma$ 与
两侧源曲线之间的最小水平距离。由 $|g_m|\leq\ee^{-|\gamma_m|\delta_\Gamma}/(2\sqrt d|\gamma_m|)$
与 $|\partial_{\nu'}g_m|\leq(|\gamma_m|^2+\beta_m^2)^{1/2}|g_m|$，每侧的振幅满足
\begin{align}
  |\phi_m|
  &\leq\Phi_m
  :=\frac{\ee^{-|\gamma_m|\delta_\Gamma}}{2\sqrt d\,|\gamma_m|}
  \Bigl(\lVert\omega\rVert_{L^1}+(|\gamma_m|^2+\beta_m^2)^{1/2}\lVert\tau\rVert_{L^1}\Bigr),
  \notag\\
  \sum_{|m|>M_{\mathrm{ev}}}
  \bigl(|\jump{\partial_\nu\widehat u}_{\Gamma,m}|^2+\ell_*^{-2}|\jump{\widehat u}_{\Gamma,m}|^2\bigr)
  &\leq2\sum_{|m|>M_{\mathrm{ev}}}(|\gamma_m|^2+\ell_*^{-2})
  \bigl((\Phi_m^{(1)})^2+(\Phi_m^{(2)})^2\bigr),
  \label{eq:artificial-tail-bound}
\end{align}
上标表示两侧，$L^1$ 范数取在相应源曲线上，不超过三角系数绝对值之和乘以曲线长度。
由于 $|\gamma_m|\geq|\beta_m|-\kappa_{\mathrm e}$ 线性增长，右端各项按
$\exp(-4\pi|m|\delta_\Gamma/d)$ 乘多项式因子衰减，可用几何级数写出闭式上界。
\end{workonly}

\subsubsection{半波导的无穷求和}

由式~\eqref{eq:half-guide-raw-field} 的平移结构，$\partial\Omega_{j,\pm}$ 与
$\Gamma_{j,\pm}$ 上的跳跃只依赖状态 $\mathsf P_\pm^jc_\pm$。对 $v\in\CC^K$，定义
$b_\pm(v)$ 为两项之和：$u_\pm[v]$ 在 $\partial\Omega_\pm$ 上的跳跃积分，以及共享端口
$\Gamma_{0,\pm}$ 上的跳跃积分（侧 $1$ 为 $u_\pm[v]$，侧 $2$ 为平移一个周期的
$u_\pm[\mathsf P_\pm v]$），均按式~\eqref{eq:boundary-defect-norm} 中各界面所属组的权重计量。记 $b_0$
为 $\partial\Omega_0$（若存在）及 $\Gamma_-$、$\Gamma_+$ 上的跳跃积分之和，后两者分别
比较中心场与 $u_-[c_-]$、$u_+[c_+]$。于是
\begin{align}
  \sum_g\varepsilon_g^2
  &=b_0+\sum_{\pm}\sum_{j=0}^{\infty}b_\pm(\mathsf P_\pm^jc_\pm).
  \label{eq:full-propagation-boundary-sum}
\end{align}
每条共享端口归属于较靠近中心的胞元，没有重复计数。全部跳跃线性依赖于状态，故
$b_\pm(v)=\lVert\mathcal L_\pm v\rVert^2$，$\mathcal L_\pm$ 是从 $\CC^K$ 到相应加权
$L^2$ 空间的线性映射；计算中不需要组装 Gram 矩阵。


\pending{2026-11-06：尾项常数 $\Theta_\pm$ 改为 $\lVert\mathcal L_\pm\rVert^2$ 的任意上界；原 Hilbert--Schmidt 取法只在工作版保留}
有限部分从 $v=c_\pm$ 开始，重复计算 $b_\pm(v)$ 并更新 $v\leftarrow\mathsf P_\pm v$，共
$N_{\mathrm{cell}}$ 次；只对实际出现的状态计算，不必建立全部端口列的映射。尾项不能仅以
末项很小为由丢弃。设常数 $\Theta_\pm$ 满足
\begin{align}
  b_\pm(v)&\leq\Theta_\pm\lVert v\rVert_2^2\qquad(v\in\CC^K),
  \label{eq:lead-operator-bound}
\end{align}
即 $\Theta_\pm\geq\lVert\mathcal L_\pm\rVert^2$。设已验证某正整数 $p$ 使
$\lVert\mathsf P_\pm^p\rVert_2\leq r_\pm<1$。把 $j=N_{\mathrm{cell}}+ip+l$
（$0\leq l<p$）代入并使用
$\lVert\mathsf P^{ip+l}w\rVert_2\leq\lVert\mathsf P^l\rVert_2r^i\lVert w\rVert_2$，得
\begin{align}
  \sum_{j=N_{\mathrm{cell}}}^{\infty}b_\pm(\mathsf P_\pm^jc_\pm)
  &\leq\frac{\Theta_\pm\sum_{l=0}^{p-1}\lVert\mathsf P_\pm^l\rVert_2^2}{1-r_\pm^2}
  \bigl\lVert\mathsf P_\pm^{N_{\mathrm{cell}}}c_\pm\bigr\rVert_2^2.
  \label{eq:boundary-propagation-tail}
\end{align}
若 $\lVert\mathsf P_\pm\rVert_2<1$，取 $p=1$。$\mathsf P_\pm$ 一般不正规，即使谱半径
小于一，$\lVert\mathsf P_\pm\rVert_2$ 也可能大于一，此时须取 $p>1$。整个计算保留完整的
$\mathsf P_\pm$，不假定传播方向正交或逐模独立。本文的 $\Theta_\pm$ 由解析包络给出
（点源行范数、近像层势与自作用对数核的上界，附录~\ref{app:reliable-eps}），无须对 $K$ 个
基向量逐列评价；它较粗，由略增大的 $N_{\mathrm{cell}}$ 补偿。


\begin{workonly}
有限部分从 $v=c_\pm$ 开始，重复计算 $b_\pm(v)$ 并更新 $v\leftarrow\mathsf P_\pm v$，共
$N_{\mathrm{cell}}$ 次。尾项不能仅以末项很小为由丢弃。令 $\boldsymbol e_s$ 为 $\CC^K$
的标准基，则
\begin{align*}
  b_\pm(v)\leq\lVert\mathcal L_\pm\rVert_{\mathrm{HS}}^2\lVert v\rVert_2^2
  =\Theta_\pm\lVert v\rVert_2^2,
  \qquad\Theta_\pm:=\sum_{s=1}^Kb_\pm(\boldsymbol e_s),
\end{align*}
$\Theta_\pm$ 只需对 $K$ 个基向量各调用一次同一评价过程。设已验证某正整数 $p$ 使
$\lVert\mathsf P_\pm^p\rVert_2\leq r_\pm<1$。把 $j=N_{\mathrm{cell}}+ip+l$
（$0\leq l<p$）代入并使用
$\lVert\mathsf P^{ip+l}w\rVert_2\leq\lVert\mathsf P^l\rVert_2r^i\lVert w\rVert_2$，得
\begin{align}
  \sum_{j=N_{\mathrm{cell}}}^{\infty}b_\pm(\mathsf P_\pm^jc_\pm)
  &\leq\frac{\Theta_\pm\sum_{l=0}^{p-1}\lVert\mathsf P_\pm^l\rVert_2^2}{1-r_\pm^2}
  \bigl\lVert\mathsf P_\pm^{N_{\mathrm{cell}}}c_\pm\bigr\rVert_2^2.
  \label{eq:boundary-propagation-tail-old}
\end{align}
若 $\lVert\mathsf P_\pm\rVert_2<1$，取 $p=1$。$\mathsf P_\pm$ 一般不正规，即使谱半径
小于一，$\lVert\mathsf P_\pm\rVert_2$ 也可能大于一，此时须取 $p>1$。整个计算保留完整的
$\mathsf P_\pm$，不假定传播方向正交或逐模独立。

\pending{尾项常数的另一取法（2026-11-05 新增）}式~\eqref{eq:boundary-propagation-tail-old} 只用到
$b_\pm(v)\leq\Theta_\pm\lVert v\rVert_2^2$，$\Theta_\pm$ 可换成 $\lVert\mathcal L_\pm\rVert^2$
的任意上界。按定义计算 $\Theta_\pm$ 要对全部 $K$ 个基向量各算一遍跳跃，每个目标都要带
$K$ 列（左右端口合计时为 $2K$ 列），在实现中是主要开销。本文改用解析包络：点源行范数、
近像层势与自作用对数核的上界，给出较粗但无需逐列评价的常数 $B_\pm\geq\lVert\mathcal L_\pm\rVert$。
有限部分只对实际出现的状态 $\mathsf P_\pm^jc_\pm$（$j<N_{\mathrm{cell}}$）计算，由线性性
与逐列计算的结果相同；尾项较粗，由略增大的 $N_{\mathrm{cell}}$ 补偿。
\end{workonly}

\subsubsection{场范数}
\label{subsec:mass}

同样由平移结构，
\begin{align}
  U^2&=\int_{\mathcal C_0}\rho|\widehat u|^2
    +\sum_{\pm}\sum_{j=0}^{\infty}m_\pm(\mathsf P_\pm^jc_\pm),
  &m_\pm(v)&:=\int_{\mathcal C_\pm}\rho|u_\pm[v]|^2.
  \label{eq:mass-propagation-sum}
\end{align}
把 $b_\pm$ 换成 $m_\pm$，式~\eqref{eq:boundary-propagation-tail} 的论证同样适用，这就
核验了 (P4)。定理只需要 $U$ 的下界：对任意可测子集 $D\subset B$，
\begin{align}
  \Bigl(\int_D\rho|\widehat u|^2\Bigr)^{1/2}&\leq U.
  \label{eq:mass-lower-bound}
\end{align}
\pending{圆盘下界（2026-11-04 新增）}最便宜的选择是均匀分片内的圆盘。设闭圆盘
$\overline{D_r(\mathbf c)}$ 含于系数为常数 $\rho_b$ 的一个分片内部，
$\kappa_b=\sqrt{\mu\rho_b}$。圆周平均 $m(s)=\frac1{2\pi}\int_0^{2\pi}\widehat u(\mathbf c+s\ee^{\ii\phi})\,\dd\phi$
满足径向 Helmholtz 方程且 $m(0)=\widehat u(\mathbf c)$，故 $m(s)=J_0(\kappa_bs)\widehat u(\mathbf c)$。
由 Cauchy--Schwarz，圆周上的平方积分至少为 $2\pi s\lvert m(s)\rvert^2$，对 $s$ 积分得
\begin{align}
  \int_{D_r(\mathbf c)}\rho|\widehat u|^2
  \geq\rho_b\pi r^2\bigl[J_0(\kappa_br)^2+J_1(\kappa_br)^2\bigr]\,\lvert\widehat u(\mathbf c)\rvert^2 .
  \label{eq:disk-mass-lower-bound}
\end{align}
互不相交的若干圆盘可以相加；也可用圆盘上的 Fourier--Bessel 系数得到更强的下界。
可选 $D$ 为中心胞元或若干相邻胞元中去掉界面细邻域后的区域。在这样的 $D$ 上层势被积
函数光滑，用原始 $N_\eta$ 个节点的梯形公式即可准确评价场值，不需要近奇异求积。$D$ 应
包含场的主要质量，否则下界过小会使条件~\eqref{eq:positive-denominator} 失效。

\subsubsection{可靠评价}
\label{subsec:reliable-evaluation}

\pending{2026-11-05 新增；2026-11-06 只保留原理，Chebyshev 误差界与节点值的三部分移入附录~\ref{app:reliable-eps}}
上面的周期梯形公式只给估计。可靠上界按目标面板计算。把一条界面的参数区间分成面板
$I$，记加权跳跃向量
\begin{align*}
  f(t):=|\boldsymbol r'(t)|^{1/2}\bigl(\jump{\partial_\nu\widehat u}_\Gamma,\ 
  \ell_g^{-1}\jump{\widehat u}_\Gamma\bigr)(\boldsymbol r(t)).
\end{align*}
若多项式向量 $P_I$ 满足 $\sup_{t\in I}|f(t)-P_I(t)|\leq e_I$，则
\begin{align}
  \lVert f\rVert_{L^2(I)}
  &\leq\sqrt{V_I}+\sqrt{|I|}\,e_I,
  &V_I&:=\int_I|P_I|^2\,\dd t,
  \label{eq:panel-defect-bound}
\end{align}
$V_I$ 是多项式的精确积分。各面板的平方上界相加再开方，就是该界面的可靠上界。$P_I$ 取
Chebyshev 插值，$e_I$ 由解析延拓上的包络与节点值误差给出；节点值与 $V_I$ 都用球算术
计算。材料界面上两侧的场各约为 $1$，跳跃约为 $10^{-12}$，约十二位有效数字在相减中
消去；保留 $P_I$ 各部分的符号再积分，球半径只需覆盖相减之后的量。细节见
附录~\ref{app:reliable-eps}。



\subsubsection{计算流程}

\begin{enumerate}
  \item 由式~\eqref{eq:candidate-vector} 得到 $(\xi,c_-,c_+)$，并固定其倍数；拟合胞元核
    系数 $\alpha$，从而固定场 $\widehat u$；
  \item 对每组计算中心胞元及两侧 $j<N_{\mathrm{cell}}$ 个胞元的跳跃积分，加上尾项
    ~\eqref{eq:boundary-propagation-tail}，开平方得 $\varepsilon_g$；
  \item 由式~\eqref{eq:mass-lower-bound} 与圆盘下界~\eqref{eq:disk-mass-lower-bound}
    计算 $U$ 的下界；
  \item 由式~\eqref{eq:joint-multiplier-constant} 计算 $C_g$，按
    式~\eqref{eq:trace-defect-product} 置 $Q$，检查 $Q<U$，代入
    式~\eqref{eq:spectral-distance-bound}。
\end{enumerate}
\pending{2026-11-06：流程改为分组形式，删去 $\Theta_\pm$ 的逐列计算}
作数值估计时，逐步增加评价参数，直到界面积分与传播尾项的变化都低于预先选定的容差。
这样得到的是谱距离的数值估计；加密稳定性不能代替可靠误差界。

% ====================================================================
\subsection{可靠上界与谱的解释}
\label{subsec:certified}

\pending{2026-11-06：改为分组输入并附输入一览；原稿“舍入可以只给估计”的退路删去，本文全部输入均计入舍入}
若要得到可靠上界，须为定理的输入建立单侧界：
\begin{align}
  C_g&\leq\overline C_g,
  &\varepsilon_g&\leq\overline\varepsilon_g,
  &0<\underline U&\leq U,
  &\underline D&:=\underline U-\sum_g\overline C_g\,\overline\varepsilon_g>0.
  \label{eq:certified-inputs}
\end{align}
各量必须来自同一候选、同一倍数与同一组密度定义的同一个场。表~\ref{tab:certified-inputs}
列出每个输入的计算方法；全部计算用球算术或向外舍入，舍入误差都已计入
（附录~\ref{app:reliable-ball}）。

\begin{table}[htbp]
  \centering
  \caption{可靠上界的输入。}
  \label{tab:certified-inputs}
  \begin{tabular}{lll}
    \toprule
    输入 & 内容 & 计算 \\
    \midrule
    $\overline\varepsilon_g$ & 面板 Chebyshev 界、自作用、光滑展开、局部项、半波导尾 & 附录~\ref{app:reliable-eps} \\
    $\underline U$ & 若干圆盘上的质量下界，只需圆心点值 & 附录~\ref{app:reliable-mass} \\
    $\overline C_g$ & 领子几何条件的验证与联合谱乘子的上确界 & 附录~\ref{app:reliable-cg} \\
    $\overline a_s,\underline b_s$ & 平面波矩阵的惯性测试（第~\ref{subsec:gap-certificate}~小节） & 附录~\ref{app:reliable-band} \\
    $R_{\mathrm{spec}}$，区间 & 标量装配，端点向外取整 & 附录~\ref{app:reliable-assembly} \\
    \bottomrule
  \end{tabular}
\end{table}

由定理~\ref{thm:spectral-distance} 立即得到
\begin{align}
  R_{\mathrm{spec}}
  &:=\frac{2\mu\sum_g\overline C_g\,\overline\varepsilon_g}{\underline D},
  &\operatorname{dist}\bigl(\mu,\sigma(A(\beta))\bigr)&\leq R_{\mathrm{spec}}.
  \label{eq:certified-spectral-distance}
\end{align}
若 $\underline D\leq0$，公式不给出有限的界。$R_{\mathrm{spec}}$ 是到谱的距离上界，不是
残量范数。因此向外舍入的区间
\begin{align}
  I_{\mathrm{spec}}&:=[\mu-R_{\mathrm{spec}},\mu+R_{\mathrm{spec}}],
  &I_k&:=\bigl[\sqrt{\mu-R_{\mathrm{spec}}},\sqrt{\mu+R_{\mathrm{spec}}}\bigr]
  \label{eq:certified-spectral-interval}
\end{align}
必与 $\sigma(A(\beta))$ 相交（$I_k$ 为对应的波数区间）。这不识别最近谱点的编号或重数，
也不保证区间内只有一个谱点。若另行证明 $I_{\mathrm{spec}}$ 含于某个与
$\sigma_{\mathrm{ess}}(A(\beta))$ 不交的开区间（第~\ref{subsec:gap-certificate}~小节的
带隙证书~\eqref{eq:gap-condition}），则由命题~\ref{prop:discrete-spectrum-in-gap}，
$I_{\mathrm{spec}}$ 中至少有一个孤立、有限重的特征值，即一个导模。

\begin{workonly}
原稿对 $\overline\varepsilon_\partial$ 的要求中有一句：“舍入可以只给估计，此时结论是‘截断误差意义下可靠’”。本文的实现全部计入舍入，这一退路不再需要。
\end{workonly}

